\subsection{Условия}
Ветвления выполняются через \verb|if ... then ... elseif ... else ... end|. Истинными считаются любые значения, кроме \verb|false| и \verb|nil|. Шаблоны: «охранные условия» (ранний выход), таблица переходов (эмуляция конвейера состояний), каскад \verb|elseif| для диапазонов.
\begin{minted}{lua}
-- классификация числа x
if x > 0 then
  state = "positive"
elseif x == 0 then
  state = "zero"
else
  state = "negative"
end
\end{minted}
\begin{minted}{lua}
-- безопасно работаем с полями, если таблица и поле существуют
if t and t.value then
  result = t.value * 2
end
\end{minted}
\paragraph{Пояснения.}
- «Охранное» условие позволяет быстро завершить ветку при неверных входных данных.
- Вложенные условия стоит разбивать на логические блоки для читаемости.
- Для выбора из фиксированного набора значений удобно использовать таблицу функций: \verb|actions[state]()|.
\subsection*{Задачи для самостоятельного решения}
\begin{enumerate}
  \item Напишите условие, присваивающее \verb|sign="+"| если \verb|x>0|, иначе \verb|sign="-"|.
  \item Присвойте \verb|parity="even"| при чётном \verb|n|, иначе \verb|"odd"|.
  \item Составьте ветвление: \verb|temp<0| → \verb|"cold"|, \verb|0<=temp<20| → \verb|"cool"|, иначе \verb|"warm"|.
  \item Присвойте \verb|level="low"| если \verb|v<3.7|, \verb|"ok"| если \verb|v>=3.7|.
  \item Напишите условие: если \verb|config.enabled| истинно, то \verb|mode="run"|, иначе \verb|"stop"|.
  \item Реализуйте проверку: если \verb|t| не \verb|nil| и \verb|t.x| существует, удвойте \verb|t.x|.
  \item Напишите ветвление, присваивающее цвет: \verb|r=1,g=0,b=0| при \verb|state=="RED"|, \verb|r=1,g=1,b=0| при \verb|"YELLOW"|, иначе \verb|r=0,g=1,b=0|.
  \item Составьте условие: если \verb|s==""|, то \verb|empty=true|, иначе \verb|false|.
  \item Реализуйте защиту от деления на ноль: если \verb|b~=0|, то \verb|q=a/b|.
  \item Напишите проверку «строка \verb|name| начинается с \verb|"P"|» (используйте \verb|string.sub|).
  \item Если \verb|voltage<3.5|, установите \verb|alarm=true|.
  \item Составьте ветвление: если \verb|t.mode=="manual"| и \verb|t.enabled|, то \verb|run=true|.
  \item Если \verb|count==0|, присвойте \verb|count=1|, иначе увеличьте на \verb|1|.
  \item Напишите условие, которое инвертирует булеву переменную \verb|on|.
  \item Реализуйте тройное ветвление по \verb|speed|: \verb|slow|, \verb|medium|, \verb|fast|.
\end{enumerate}
